\subsection{Teorema de Sprague-Grundy e Nim}

O teorema descreve o jogo de dois players inicial i.e. aqueles em que os movimentos disponíveis e winning/losing depende apenas do estado do jogo. Ou seja, a única diferença entre os jogadores é quem se move primeiro.

Também assumimos que o jogo tem informação perfeita i.e. os jogadores sabem perfeitamente as regras e movimentos possíveis. Também assumimos que é finito i.e. depois de uma certa quantidade de movimentos um dos jogadores acabará numa losing position da qual não pode se mover. Então não há empate nesse jogo.

Esses jogos podem ser descritos completamente por um DAG.

\subsubsection{Nim}

Esse é um jogo que obedece tudo o que foi dito acima. Além disso, qualquer jogo imparcial two-layer de informação perfeita pode ser reduzido para o jogo de nim. Então estudar o nim resolve todos os jogos.

Você tem várias pilhas com várias pedras. Num movimento, um jogador deve escolher uma pilha e retirar uma quantidade inteira positiva de pedras dessa pilha. Ele perde se não tem mais pedras.

\textbf{Teorema (Nim):} O jogador atual tem uma estratégia vencedora se e só se a xor-sum das pilhas é não zero i.e. $s = a_1 \oplus \dots \oplus a_n \neq 0$. Além disso, o movimento vencedor de um dado estado consiste em reduzir alguma pilha de tamanho $x$ para o tamanho $s \oplus x$.

\textbf{Corolário:} todo estado do nim pode ser substituído por outro com o mesmo xor. Quando estamos analisando um nim com várias pilhas podemos observar isso como substituindo por uma única pilha de tamanho $s$.

\subsubsection{Teorema de Sprague-Grundy}

Agora suponha que podemos adicionar pedras a uma pilha escolhida. As regras exatas de como e quando pode ocorrer esse aumento não nos interessam, mas o que importa é que essas regras devem deixar nosso jogo acíclico.

\textbf{Lema:} A adição de acréscimos a pedras não muda a maneira como estados de vitória/derrota são determinados. Ou seja, podemos ignorá-los.

\textbf{Teorema:} Considere um estado $v$ de um jogo imparcial e sejam $v_i$ os estados alcançáveis dele (com $i \in \{1, \dots, k\}$, $k \geq 0$). A esse estado podemos designar um jogo de nim completamente equivalente com uma pilha de tamanho $x$. O número $x$ é dito o valor de Grundy ou nim-value do estado $v$. Esse número pode ser achado de maneira recursiva: 
$$x = \text{mex} \{ x_1 , \dots, x_k\}$$
em que $x_i$ são os valores de grundy para os estados $v_i$ e o mex pega o menor valor não negativo que não está no conjunto.

\textbf{Aplicação:} Para calcular o valor de grundy de um estado, você deve
\begin{enumerate}
    \item Pegar todas as transições a partir desse estado
    \item Cada transição pode levar à soma de jogos independentes. Calcule o valor de grundy para cada jogo independente e xor-sum neles.
    \item Depois de calcular o valor de grundy para cada transição, pega o mex dos números
    \item Se o valor for zero, então o estado atual é losing. Se não for, é winning.
\end{enumerate}

Dica: em tarefas específicas estudar a tabela de valores em busca de padrões é uma ótima ideia. Em muitos jogos onde uma análise teórica pode parecer difícil, os valores de grundy são periódicos ou então tem uma forma compreensível. Em 99,9\% dos casos dá pra provar o padrão por indução. Entretanto, ainda está aberto o problema de se existem regularidades em alguns casos.

\prob{1}{grundy-1} Considere uma fita de tamanho $1 \times n$. Em um movimento, o jogador deve colocar uma cruz. Mas é proibido colocar duas cruzes uma ao lado da outra. O jogador sem movimentos válidos perde.

\textbf{Solução:} Seja $g$ o valor de grundy. É trivial ver que
$$g(n) = \text{mex} \left( \{ g(n-2) \} \cup \{ g(i-2) \oplus g(n-i-1) | 2 \leq i \leq n-1\} \right)$$
Com isso, temos uma solução $O(n^2)$. Na verdade, $g(n)$ tem período de tamanho $34$ começando com $n=52$.